\section{Conclusion}

% We introduce Radial Consensus Score (RCS), a simple geometric method for best-of-$N$ selection in LLMs. RCS computes a weighted semantic center via the Fréchet mean of answer embeddings and ranks candidates by their distance to this center, capturing consensus while optionally incorporating frequency or generation probability. 
% It consistently outperforms strong baselines across diverse benchmarks, with larger gains under higher sampling budgets, while being efficient, black-box compatible, and applicable as a drop-in replacement for majority voting in multi-agent debate. This highlights geometric consensus as a simple and scalable principle for reliable answer selection.

We introduce Radial Consensus Score (RCS), a geometric framework for best-of-$N$ answer selection in LLMs. RCS ranks candidates by their distance to a weighted Fréchet mean of answer embeddings, with the weighting scheme---uniform, frequency, probability, or cosine-similarity---determining how consensus is defined. Compared to Semantic Consensus Weighting (SCW), which applies the same cosine signal via a discrete grouping step, RCS$_\text{cosine}$ uses it continuously as Fréchet-mean weights, recovering higher-quality minority answers that discrete voting misses. RCS variants consistently outperform strong baselines across diverse benchmarks, with gains growing at higher sampling budgets, while remaining black-box compatible and applicable as a drop-in replacement for majority voting in multi-agent debate.

% In the special case of equivalent candidates, RCS reduces to frequency- or probability-based selection, recovering self-consistency in low-diversity settings. Overall, RCS highlights the effectiveness of geometric consensus for robust and semantically aware answer selection beyond voting-based approaches.